% 图 2-1：按真实接口分区，父框由 fit 自动包围子节点。
\begin{tikzpicture}[font=\footnotesize,
 box/.style={draw=softline,rounded corners=2pt,align=center,inner sep=5pt,text width=2.8cm,minimum height=1.05cm,fill=white},
 frame/.style={draw=macline,rounded corners=3pt,inner sep=6mm},
 ar/.style={-{Stealth[length=1.5mm]},thick,draw=softline},
 ctl/.style={ar,dashed},lab/.style={font=\scriptsize,fill=white,inner sep=2pt}]
\node[font=\small\bfseries,anchor=west] (title) at (-7.3,1.2) {当前演示：pynq\_puf3\_wrapper};
\node[box,fill=macpeach] (ps) at (-5.7,0) {PS7\\FCLK0 100 MHz};
\node[box,fill=macpeach] (win) at (-1.9,0) {AXI-Lite\\CONTROL / MAILBOX};
\node[box,fill=macgray] (rst) at (1.9,0) {复位同步\\SOC\_RESET};
\node[box,fill=macblue] (bram) at (5.7,0) {BRAM 128 KiB\\32 bit $\times$ 32768};
\draw[ar] (ps.east)--(win.west);
\node[font=\small\bfseries,anchor=west] (stitle) at (-7.3,-1.9) {片上系统：soc\_top};
\node[box,text width=6.4cm,fill=macgreen,minimum height=1.65cm] (cpu) at (-3.7,-3.2) {\textbf{cpu\_core：五级流水线}\\IF / ID / EX / MEM / WB\\I-cache、D-cache 各 2 KiB（两路）};
\node[box,text width=5.3cm,fill=maclav,minimum height=1.65cm] (control) at (3.8,-3.2) {\textbf{CPU 内部控制与反馈}\\64 项 BTB / BHT\\旁路 bypass；冒险控制 stall\_ctrl};
\node[box,text width=6.4cm,fill=macamber] (ise) at (-3.7,-4.85) {\textbf{新增 Ascon ISE（位于 EX）}\\custom-0；320-bit 状态；P8 / P12};
\draw[ar] (cpu.south)--(ise.north);
\draw[ctl,<->] (cpu.east)--(control.west);
\node[box,text width=13.6cm,fill=macgreen,minimum height=1.1cm] (bus) at (0,-6.55) {bus\_interconnect：内部地址译码与仲裁\\BRAM 访问经 cpu\_axi\_bridge $\rightarrow$ AXI4 master};
\draw[ar] (cpu.west)--(-7.3,-3.2)--(-7.3,-6.55)--(bus.west);
\node[box,text width=2.7cm,fill=macblue] (uart) at (-5.7,-8.35) {UART\\0x2000\_0000};
\node[box,text width=2.7cm,fill=macblue] (gpio) at (-1.9,-8.35) {GPIO\\0x2001\_0000};
\node[box,text width=2.7cm,fill=macblue] (clint) at (1.9,-8.35) {CLINT\\0x1200\_0000};
\node[box,text width=2.7cm,fill=macblue] (cpi) at (5.7,-8.35) {CPI 事件计数\\0x6000\_0000};
\foreach \n in {uart,gpio,clint,cpi}{\draw[ar,<->] (bus.south -| \n.north)--(\n.north);}
\node[box,text width=9.5cm,fill=macamber] (asc) at (0,-10.05) {\textbf{Ascon MMIO：0x3000\_0000}\\调用者 key / RO-PUF + DNA 派生 key（HWKEY 显式选择）\\密钥控制窗口 +0x80；CPU 搬运数据并轮询完成};
\draw[ar,<->] (bus.east)--(7.55,-6.55)--(7.55,-10.05)--(asc.east);
\draw[ar,<->] (bram.east)--(7.9,0)--(7.9,-6.55)--(bus.east);
\node[lab,rotate=90] at (7.9,-3.0) {AXI / BRAM 数据通路};
\draw[ctl] (win.south)--(-1.9,-1.15)--(-7.6,-1.15)--(-7.6,-2.35)--(cpu.north west);
\begin{scope}[on background layer]
 \coordinate (routebound) at (8.2,-6.55);
 \node[frame,inner sep=10mm,fill=macpeach!25,fit=(title)(ps)(win)(rst)(bram)(stitle)(cpu)(control)(ise)(bus)(uart)(gpio)(clint)(cpi)(asc)(routebound)] {};
 \node[frame,fill=macblue!18,fit=(stitle)(cpu)(control)(ise)(bus)(uart)(gpio)(clint)(cpi)(asc)] (socframe) {};
\end{scope}
\draw[ctl] (rst.south)--(socframe.north -| rst.south);
\end{tikzpicture}
